\section{Special rays and the end of the combination proof}
\label{ray:section}

We finish the proof of Theorem~\ref{ann:combination}. Suppose that $G$
is not hyperbolic. Propositions~\ref{ann:trajectory}
and~\ref{stat:stationary} then provide a
nonzero stationary current $\mu$, its support $\Omega$, and a conull
invariant set $\mathcal U$. By Corollary~\ref{exp:support}, every endpoint
of every line of $\Omega$ is irrational, and every support line has a unique bilateral trajectory. Each fixed
finite-path transfer maps its domain in $\Omega$ into $\Omega$. Moreover, the cores
$\Gamma_n$ have a fixed small subgraph along the chosen subsequence, and
there is a uniform bound on the length of a segment of an $\Omega$-line
that lifts entirely to that small subgraph.

The remaining argument detects branching in the language of $\Omega$.
Bounded natural graph complexity leaves only finitely many rays with
more than one possible past. A periodic orbit among their endpoints
gives a partial conjugation moving points progressively farther along
one ray. Its invariant measure cannot be finite on the lines through
a vertex.

\subsection{Finitely many rays with two pasts}

Fix the basis of $P$ used above, and write $T_P$ for its Cayley tree.
Orient each line from its negative endpoint to its positive endpoint.
An infinite reduced word is \emph{occurring} if it labels a positive ray
in a line of $\Omega$. An occurring word $\rho$ is \emph{left-special}
if there are distinct basis letters $a,b$ such that both $a\rho$ and
$b\rho$ occur. These extensions are required to be reduced. Translating
their rays to start at $1$ gives two lines of $\Omega$ with a common
positive ray and different incoming edges.

\begin{lemma}\label{ray:finite}
The set of left-special occurring words is finite and nonempty.
\end{lemma}
\begin{proof}
Let $\rho$ be left-special. Both of its predecessor lines lift to each
$\Gamma_n$ by Lemma~\ref{car:window-cores}. Their common endpoint is
irrational, so the uniqueness of a starting state for an irrational
ray puts both lifts at the same vertex. This vertex has two different
incoming edges as well as the outgoing ray edge. It is therefore a
natural vertex.

Let $B$ bound the lengths of small-only segments of support lines, as in
Corollary~\ref{exp:support}. Starting at this natural vertex, the lift
of $\rho$ follows at most $B$ labelled edges in the small subgraph before
entering a long natural edge. There are only finitely many possible
prefix words of length at most $B$ and finitely many oriented natural
edge types. On a subsequence the prefix and the oriented type are fixed.
The lengths of these long edges tend to infinity, and their oriented
labels have a fixed infinite prefix limit by the construction of
$\Gamma_n$. Thus $\rho$ is the concatenation of one of finitely many
prefixes and one of finitely many prescribed infinite words. This proves
finiteness.

Suppose now that no occurring ray is left-special. Then there is an
integer $N$ such that every occurring word of length $N$ has a unique
predecessor letter. Indeed, otherwise arbitrarily long finite words
would have two distinct predecessor letters. Passing to a fixed pair
of letters and a convergent subsequence of their full-line extensions
would give a left-special infinite word. This uses compactness of the
space of lines through a vertex and the closedness of $\Omega$.

Form the finite directed graph whose vertices are occurring words of
length $N$ and whose edges are occurring words of length $N+1$, an edge
running from its length-$N$ prefix to its length-$N$ suffix. Every vertex
has indegree one and outdegree at least one. The sum of the indegrees is
the number of vertices, so every outdegree is also one. The graph is a
disjoint union of directed cycles. Every line in $\Omega$ would therefore
have a periodic label, contrary to Proposition~\ref{exp:no-periodic}.
\end{proof}

\subsection{A return map with positive lag}

Let $\mathcal E$ be the set of $P$-orbits of the positive endpoints of
left-special rays. This is a finite nonempty set. The forward path
transfer induces a permutation of $\mathcal E$. To see this, take two
predecessor lines of such a ray. Their common irrational endpoint puts
them in the same unique domain coset. The boundary embedding on that
coset sends them to two support lines with the same positive endpoint
and distinct negative endpoints. After their positive rays merge, the
two incoming edges exhibit a left-special ray. A change of the input
representative by an element of $P$ only changes the output by an
element of $P$. The backward transfer is the inverse construction, so
the resulting map on $\mathcal E$ is a permutation.

Choose a left-special word based at $1$ whose endpoint $\xi$ represents
a periodic element of $\mathcal E$, of period $p>0$.
Compose the $p$ forward transfers and adjust the final frame by an
element of $P$. The actual path construction gives an element $v\in G$
with nonzero height $p$ and a partial conjugation
\begin{equation}\label{ray:return}
 \Theta:S\longrightarrow T,\qquad
 \Theta(x)=v^{-1}xv,\qquad
 S=P\cap vPv^{-1},\quad T=P\cap v^{-1}Pv,
\end{equation}
whose boundary map fixes $\xi$. Here $S$ and $T$ are the full
stabilizers, expressed in their initial and final frames, of the
corresponding finite path. In particular they are finitely generated
subgroups of $P$. The map and its inverse preserve $\mu$ on their
respective domains in $\mathcal U$.

Write $r_0=1,r_1,r_2,\ldots$ for the consecutive vertices of the ray to
$\xi$ in $T_P$. Since $\xi\in\Lambda S$ and $S$ is finitely generated,
choose $s_n\in S$ and a fixed $K\geq0$ such that
$d(s_n,r_n)\leq K$ for every $n$.
The embedding $\Theta$ is a quasi-isometry onto $T$ and fixes the
boundary endpoint $\xi$. Consequently there is a sequence of
nonnegative integers $f(n)$ with
\begin{equation}\label{ray:indices}
 d\bigl(\Theta(s_n),r_{f(n)}\bigr)\leq C,
 \qquad f(n)\longrightarrow\infty,
 \qquad |f(n+1)-f(n)|\leq C,
\end{equation}
after increasing a fixed constant $C$. These are the ray versions of
the transport estimates in Lemma~\ref{exp:transport}.

\begin{lemma}\label{ray:lag}
After replacing $\Theta$ by its inverse if necessary, the choices above
may be made so that
\[
 f(n)-n\longrightarrow+\infty.
\]
\end{lemma}
\begin{proof}
If $|f(n)-n|$ were bounded along an infinite subsequence, then
$s_n^{-1}\Theta(s_n)$ takes only finitely many values on that subsequence.
The elements $s_n$ escape every finite ball. Hence there are distinct
$s_n,s_m$ and $d\in P$ with
$\Theta(s_n)=s_nd$ and $\Theta(s_m)=s_md$. It follows that
\[
 \Theta(s_ns_m^{-1})=(s_nd)(s_md)^{-1}=s_ns_m^{-1}.
\]
The nontrivial element $s_ns_m^{-1}$ has infinite order and commutes
with $v$. Their heights are respectively zero and $p\ne0$, so they
generate a subgroup isomorphic to $\mathbb Z^2$. This contradicts the
Baumslag--Solitar exclusion.

Thus $|f(n)-n|\to\infty$. The signed sequence $f(n)-n$ has bounded
successive increments by \eqref{ray:indices}, so it is eventually
either positive or negative. The positive case is the assertion.
In the negative case choose $t_m\in T$ uniformly close to $r_m$
and integers $g(m)$ such that
$\Theta^{-1}(t_m)$ is uniformly close to $r_{g(m)}$. Because $f$ tends
to infinity with bounded increments, a first-crossing choice $n(m)$
satisfies
\[
 |f(n(m))-m|\leq C,\qquad n(m)\longrightarrow\infty.
\]
Then $t_m$ and $\Theta(s_{n(m)})$ are uniformly close. Applying the
inverse quasi-isometry gives
\[
 g(m)=n(m)+O(1),\qquad
 g(m)-m=n(m)-f(n(m))+O(1)\longrightarrow+\infty.
\]
This proves the assertion for the inverse map.
\end{proof}

\subsection{A cylinder that cannot lose mass}

After the possible inversion, use $\Theta,S,T,f$ for the resulting
map and ray data, and $K,C$ for their uniform approximation bounds.
The labels of distinct suffixes of the ray $(r_n)$ are distinct: an
equality would make the ray eventually periodic and $\xi$ rational.
By Lemma~\ref{ray:finite}, there is $n_0$ such that no suffix beginning
after $r_{n_0}$ is left-special.

Fix $s>n_0$. Any support line containing the entire directed suffix
$[r_s,\xi)$ has its preceding edge uniquely determined, then the edge
before that, and so on back to $r_{n_0}$. Moreover, its actual path for
the fixed window defining $\Theta$ is the same as the path for $\xi$:
an irrational endpoint determines that path uniquely. These two
statements have the following finite form.

\begin{lemma}\label{ray:forcing}
For each fixed $s>n_0$, there is $N_s$ such that, whenever $N\geq N_s$,
every line of $\Omega$ containing the directed segment $[r_s,r_N]$
contains $r_{n_0}$ and belongs to the domain of $\Theta$.
\end{lemma}
\begin{proof}
For the first assertion, a contrary sequence has a convergent
subsequence in the compact space of directed lines through $r_s$.
The limit belongs to $\Omega$ and contains the whole suffix
$[r_s,\xi)$. The condition of not containing the fixed finite segment
$[r_{n_0},r_s]$ is closed in this compact space. The limit would therefore
contradict the unique-past property just proved.

For the domain assertion, each competing line has an actual path for
the fixed window. Its stabilizer hull contains that line, hence meets
the vertex $r_s$. There are only finitely many such actual paths by
Lemma~\ref{stat:window-finiteness}. A contrary sequence therefore has a
subsequence with one fixed competing path. Its stabilizer limit set
contains the limiting positive endpoint $\xi$. Uniqueness for an
irrational endpoint identifies this path with the path defining
$\Theta$, which is a contradiction.
\end{proof}

For $N\geq N_s$, let
\[
 U_N=\{\ell\in\mathcal U:
              \ell\text{ contains the directed segment }[r_s,r_N]\}.
\]
This set has finite positive $\mu$-measure: the cylinder is compact and
open, meets the support, and $\mathcal U$ is conull. Its lines contain
$r_{n_0}$ and lie in the domain of $\Theta$ by Lemma~\ref{ray:forcing}.
The elements $s_{n_0}$ and $s_N$ lie within $K$ of these lines.
Lemma~\ref{exp:transport}, applied to the fixed domain inclusion
$S\hookrightarrow P$ and the fixed injection $\Theta:S\to P$,
gives a constant $E$ such that their images lie within $E$ of every
line of $\Theta(U_N)$. Equation~\eqref{ray:indices} then puts
$r_{f(n_0)}$ and $r_{f(N)}$ within $E+C$ of those image lines.
The input projections are in their positive order with separation at
least $N-n_0-2K$. For sufficiently large $N$, the order and separation
assertions of Lemma~\ref{exp:transport} put the two image projections,
and then the two ray-vertex projections, in the same positive order.

Choose a fixed integer $D>E+C$. The tree observation preceding
Lemma~\ref{exp:trimming} shows that every line of $\Theta(U_N)$
contains the directed central segment
$[r_{f(n_0)+D},r_{f(N)-D}]$ for all sufficiently large $N$.
Both $E$ and $D$ depend only on the fixed map and the uniform
approximation bounds; neither depends on $s$ or $N$.

Now choose $s>\max\{n_0,f(n_0)+D\}$, and then take $N$ beyond
$N_s$ and the preceding transport thresholds. Lemma~\ref{ray:lag}
also gives $f(N)-D\geq N+1$ for every sufficiently large $N$.
The invariance of $\mathcal U$ under the path transfers now yields
\begin{equation}\label{ray:nesting}
 \Theta(U_N)\subseteq U_{N+1}\subseteq U_N.
\end{equation}
Because $\Theta$ preserves measure on its domain, these inclusions give
\[
 \mu(U_N)=\mu(\Theta(U_N))
       \leq\mu(U_{N+1})\leq\mu(U_N).
\]
All sufficiently late $U_N$ have the same positive finite measure.
Continuity from above shows that their intersection $U_\infty$ has that
same measure. Every line in $U_\infty$ contains the full directed ray
$[r_s,\xi)$.

For each $n\geq s$, translate $U_\infty$ by $r_n^{-1}$. These sets all
consist of lines through $1$, and they all have the same positive
measure. Their positive endpoints $r_n^{-1}\xi$ are pairwise distinct:
an equality for $n\ne m$ would give a nontrivial element of $P$ fixing
$\xi$, making it rational. The translated sets are thus pairwise
disjoint. This contradicts the finite mass of the compact set of lines
through $1$.

We have contradicted the existence of the nonzero stationary current
forced by nonhyperbolicity. Therefore $G$ is hyperbolic.
This completes the proof of Theorem~\ref{ann:combination}.
